\ELhold

\ELsection{Course Outline}{Course Outline}

\begin{ELunit}

\ELfigure{m04-course-outline}{}

\end{ELunit}

\ELhold

\ELsection{Steady Electric Currents}{Steady Electric Currents}

\begin{ELunit}

\begin{itemize}
\tightlist
\item
  In the previous chapters we dealt with electrostatic problems due to electric charges at rest.

  \begin{itemize}
  \tightlist
  \item
    In this chapter we turn to moving electric charges, which produce electric current.
  \end{itemize}
\item
  Electric currents due to the motion of free charges are divided into three kinds

  \begin{itemize}
  \tightlist
  \item
    \textbf{Conduction currents:} produced in conductors and semiconductors by the drift motion of electrons or holes.
  \item
    \textbf{Electrolytic currents:} due to the motion of positive and negative ions
  \item
    \textbf{Convection currents:} due to the motion of electrons or ions in vacuum
  \end{itemize}
\item
  In this chapter we concentrate on conduction currents, which are governed by Ohm's law.
\end{itemize}

\ELfigure{m04-map}{}

\begin{itemize}
\tightlist
\item
  What is the concept of the \textbf{volume current density}, and what is the relation for the \textbf{convection current density}?
\item
  What is the relation for the \textbf{conduction current density}?
\item
  What is the point form of \textbf{Ohm's law}?
\item
  What is the \textbf{equation of continuity}, and how is \textbf{Kirchhoff's current law} derived from it?
\item
  What is \textbf{Joule's law}, which gives the power dissipated by a current flowing in a conducting body?
\end{itemize}

\end{ELunit}

\ELhold

\ELsection{Current Density and Ohm's Law}{Current Density and Ohm’s Law}

\begin{ELunit}

\ELfigure{m04-current-element}{}

\begin{itemize}
\tightlist
\item
  Consider the steady motion of charge carriers, each with charge \ELim{q}, moving with velocity \ELim{\vect{u}} across a small surface element \ELim{\Delta s}
\item
  If \ELim{N} is the number of charge carriers per unit volume, the charge passing through the surface \ELim{\Delta s} in the time \ELim{\Delta t} is
\end{itemize}

\ELdisp{\Delta Q\ELbr =Nq\vect{u}\cdot\uvec{n}\,\Delta s\,\Delta t\qquad\ELbr (\mathrm{C})}

\begin{itemize}
\tightlist
\item
  Since current is the time rate of change of charge, we have
\end{itemize}

\ELdisp{\Delta I\ELbr =\frac{\Delta Q}{\Delta t}\ELbr =Nq\vect{u}\cdot\uvec{n}\,\Delta s\ELbr =Nq\vect{u}\cdot\Delta\vect{s}\qquad\ELbr (\mathrm{A})}

\begin{itemize}
\tightlist
\item
  We define the volume current density, in amperes per square metre, as follows
\end{itemize}

\ELdisp{\vect{J}\ELbr =Nq\vect{u}\qquad\ELbr (\mathrm{A/m^2})}

\end{ELunit}

\begin{ELimportant}{Convection current density}

\ELdisp{\vect{J}\ELbr =\rho\vect{u}\qquad\ELbr (\mathrm{A/m^2})}

\end{ELimportant}

\begin{ELjoin}

\begin{itemize}
\tightlist
\item
  Therefore
\end{itemize}

\ELdisp{\Delta I\ELbr =\vect{J}\cdot\Delta\vect{s}}

\begin{itemize}
\tightlist
\item
  Thus the total current flowing through an arbitrary surface \ELim{S} is
\end{itemize}

\begin{ELimportant}{}

\ELdisp{I\ELbr =\int_S\vect{J}\cdot d\vect{s}\qquad\ELbr (\mathrm{A})}

\end{ELimportant}

\end{ELjoin}

\begin{ELunit}

\begin{itemize}
\tightlist
\item
  Conduction currents are the result of the drift motion of charge carriers under the influence of an applied electric field.
\item
  For most conducting materials the average drift velocity is directly proportional to the electric field intensity.

  \begin{itemize}
  \tightlist
  \item
    For metallic conductors we have
  \end{itemize}
\end{itemize}

\ELdisp{\vect{u}\ELbr =-\mu_e\vect{E}\qquad\ELbr (\mathrm{m/s})}

\begin{itemize}
\tightlist
\item
  \ELim{\mu_e}: electron mobility, in \ELim{\mathrm{m^2/V\cdot s}}
\end{itemize}

\ELdisp{\vect{J}\ELbr =-\rho_e\mu_e\vect{E}}

\end{ELunit}

\begin{ELjoin}

\begin{itemize}
\tightlist
\item
  \ELim{\rho_e=-Ne}: the volume charge density of the drifting electrons, a negative quantity
\end{itemize}

\begin{ELimportant}{Point form of Ohm's law}

\ELdisp{\vect{J}\ELbr =\sigma\vect{E}\qquad\ELbr (\mathrm{A/m^2})}

\end{ELimportant}

\end{ELjoin}

\begin{ELunit}

\begin{itemize}
\item
  \ELim{\sigma=-\rho_e\mu_e}: a fundamental parameter of the medium, called the conductivity
\item
  The unit of conductivity is \ELim{\mathrm{A/V\cdot m}}, or siemens per metre. The reciprocal of conductivity is called resistivity, in \ELim{\Omega\cdot\mathrm{m}}.
\item
  Obtaining the relation between the voltage and the current of a piece of homogeneous material of conductivity \ELim{\sigma}, length \ELim{\ell} and uniform cross-section \ELim{S} from the point form of Ohm's law
\end{itemize}

\ELfigure{m04-resistor}{}

\ELdisp{\left.\begin{aligned}V_{12}=E\ell\quad&\Longrightarrow\quad E=\frac{V_{12}}{\ell}\\I=\int_S\vect{J}\cdot d\vect{s}=JS\quad&\Longrightarrow\quad J=\frac IS\end{aligned}\right\}\quad\ELbr \frac IS\ELbr =\sigma\frac{V_{12}}{\ell}}
\ELdisp{V_{12}\ELbr =\left(\frac{\ell}{\sigma S}\right)I\ELbr =RI}

\end{ELunit}

\begin{ELimportant}{}

\ELdisp{R\ELbr =\frac{\ell}{\sigma S}\qquad\ELbr (\Omega)}
\ELdisp{G\ELbr =\frac1R\ELbr =\sigma\frac S\ell\qquad\ELbr (\mathrm{S})}

\end{ELimportant}

\begin{ELunit}

\begin{itemize}
\tightlist
\item
  \ELim{G}: conductance, the reciprocal of resistance, in mhos or siemens
\end{itemize}

\end{ELunit}

\ELhold

\ELsection{Equation of Continuity and Kirchhoff's Current Law}{Equation of Continuity and Kirchhoff’s Current Law}

\begin{ELunit}

\ELfigure{m04-map-continuity}{}

\begin{itemize}
\tightlist
\item
  The principle of conservation of charge is one of the fundamental postulates of physics.

  \begin{itemize}
  \tightlist
  \item
    Electric charge can be neither created nor destroyed.
  \end{itemize}
\item
  Consider an arbitrary volume \ELim{V} bounded by the surface \ELim{S}, containing a net charge \ELim{Q}. If a net current \ELim{I} flows out of the region through the surface, the charge in the volume must decrease at a rate equal to the current. Conversely, if a net current flows into the region through the surface, the charge in the volume must increase at a rate equal to the current.
\item
  The current leaving the region is the total outward flux of the current density vector through the surface \ELim{S}
\end{itemize}

\ELdisp{I\ELbr =\oint_S\vect{J}\cdot d\vect{s}\ELbr =-\frac{dQ}{dt}\ELbr =-\frac{d}{dt}\int_V\rho\,dv}
\ELdisp{\int_V\nabla\cdot\vect{J}\,dv\ELbr =-\int_V\frac{\partial\rho}{\partial t}\,dv}

\end{ELunit}

\begin{ELjoin}

\begin{itemize}
\tightlist
\item
  Since the above equation must hold regardless of the choice of \ELim{V}
\end{itemize}

\begin{ELimportant}{Equation of continuity}

\ELdisp{\nabla\cdot\vect{J}\ELbr =-\frac{\partial\rho}{\partial t}\qquad\ELbr (\mathrm{A/m^3})}

\end{ELimportant}

\end{ELjoin}

\begin{ELjoin}

\begin{itemize}
\tightlist
\item
  If there are no time variations
\end{itemize}

\ELdisp{\nabla\cdot\vect{J}\ELbr =0}

\begin{ELimportant}{Kirchhoff's current law}

\ELdisp{\oint_S\vect{J}\cdot d\vect{s}\ELbr =0}

\end{ELimportant}

\end{ELjoin}

\begin{ELunit}

\begin{itemize}
\tightlist
\item
  We said earlier that charges placed inside a conductor move to its surface, so that at equilibrium \ELim{\rho=0} inside the conductor. We now want to prove this statement
\end{itemize}

\ELdisp{\left.\begin{aligned}\sigma\nabla\cdot\vect{E}&=-\frac{\partial\rho}{\partial t}\\\nabla\cdot\vect{E}&=\rho/\epsilon\end{aligned}\right\}\quad\ELbr \frac{\partial\rho}{\partial t}+\frac\sigma\epsilon\rho\ELbr =0\quad\ELbr \ELbr \Longrightarrow\quad\ELbr \rho\ELbr =\rho_0e^{-(\sigma/\epsilon)t}\qquad\ELbr (\mathrm{C/m^3})}

\begin{itemize}
\tightlist
\item
  \ELim{\rho_0}: initial charge density
\item
  The above equation states that the charge density at a given point decreases exponentially with time.
\item
  The initial density decays to \ELim{1/e}, or \ELim{36.8\%}, of its value in a time equal to \ELim{\tau=\epsilon/\sigma}.

  \begin{itemize}
  \tightlist
  \item
    For a good conductor such as copper, \ELim{\tau=1.52\times10^{-19}\,\mathrm{s}}.
  \item
    For a good insulator this time may be hours or days.
  \end{itemize}
\end{itemize}

\end{ELunit}

\ELhold

\ELsection{Power Dissipation and Joule's Law}{Power Dissipation and Joule’s Law}

\begin{ELunit}

\ELfigure{m04-map-joule}{}

\begin{itemize}
\tightlist
\item
  Under the influence of an electric field, the free electrons in a conductor move and collide with the atoms of the crystal lattice. Energy is thus transferred from the electric field to the vibrating atoms.
\item
  The work done by the electric field in moving a charge \ELim{q} through the distance \ELim{\Delta\ell} is
\end{itemize}

\ELdisp{\Delta W\ELbr =\vect{F}\cdot\Delta\vect{\ell}\ELbr =q\vect{E}\cdot\Delta\vect{\ell}}

\begin{itemize}
\tightlist
\item
  This work corresponds to the power
\end{itemize}

\ELdisp{p\ELbr =\lim_{\Delta t\to0}\frac{\Delta w}{\Delta t}\ELbr =q\vect{E}\cdot\vect{u}}

\begin{itemize}
\tightlist
\item
  \ELim{\vect{u}}: the velocity of the charge carriers
\end{itemize}

\ELdisp{dp\ELbr =\rho\,dv\,\vect{E}\cdot\vect{u}}
\ELdisp{dP\ELbr =\vect{E}\cdot\vect{J}\,dv}

\end{ELunit}

\begin{ELjoin}

\begin{itemize}
\tightlist
\item
  For a given volume \ELim{V} the total power converted into heat is
\end{itemize}

\begin{ELimportant}{Joule's law}

\ELdisp{P\ELbr =\int_V\vect{E}\cdot\vect{J}\,dv\qquad\ELbr (\mathrm{W})}

\end{ELimportant}

\end{ELjoin}

\begin{ELjoin}

\begin{itemize}
\tightlist
\item
  In a conductor of constant cross-section we have
\end{itemize}

\ELdisp{P\ELbr =\int_LE\,d\ell\int_SJ\,ds\ELbr =VI}

\begin{ELimportant}{}

\ELdisp{P\ELbr =I^2R\qquad\ELbr (\mathrm{W})}

\end{ELimportant}

\end{ELjoin}

\begin{ELunit}

\begin{itemize}
\tightlist
\item
  which is the familiar relation for ohmic power and gives the heat dissipated in the resistance \ELim{R}.
\end{itemize}

\end{ELunit}

\ELhold

\ELsection{Boundary Conditions for Current Density}{Boundary Conditions for Current Density}

\begin{ELunit}

\ELfigure{m04-map-boundary}{}

\begin{itemize}
\tightlist
\item
  When a current crosses the interface between two media of different conductivities obliquely, the current density vector changes in both direction and magnitude.
\item
  The governing equations for the steady current density are
\end{itemize}

\end{ELunit}

\textbf{Governing Equations for Steady Current Density}

\Needspace{8\baselineskip}
\begin{ELltr}
\begin{longtable}[c]{@{}
  >{\raggedright\arraybackslash}p{(\columnwidth - 2\tabcolsep) * \real{0.5000}}
  >{\raggedright\arraybackslash}p{(\columnwidth - 2\tabcolsep) * \real{0.5000}}@{}}
\toprule\noalign{}
\ELtablehead \begin{minipage}[b]{\linewidth}\raggedright
Differential Form
\end{minipage} & \begin{minipage}[b]{\linewidth}\raggedright
Integral Form
\end{minipage} \\\noalign{\vskip 4pt}
\midrule\noalign{}
\endhead
\bottomrule\noalign{}
\endlastfoot
\ELim{\nabla\cdot\vect{J}=0} & \ELim{\oint_S\vect{J}\cdot d\vect{s}=0} \\\noalign{\vskip 4pt}
\ELim{\nabla\times\left(\dfrac{\vect{J}}{\sigma}\right)=0} & \ELim{\oint_C\dfrac1\sigma\vect{J}\cdot\dif\vect{\ell}=0} \\\noalign{\vskip 4pt}
\end{longtable}\end{ELltr}


\begin{ELjoin}

\begin{itemize}
\tightlist
\item
  The boundary conditions are as follows
\end{itemize}

\begin{ELimportant}{}

\ELdisp{\nabla\cdot\vect{J}\ELbr =0\quad\ELbr \ELbr \Longrightarrow\quad\ELbr  J_{1n}\ELbr =J_{2n}\qquad\ELbr (\mathrm{A/m^2})}
\ELdisp{\nabla\times(\vect{J}/\sigma)\ELbr =0\quad\ELbr \ELbr \Longrightarrow\quad\ELbr \frac{J_{1t}}{J_{2t}}\ELbr =\frac{\sigma_1}{\sigma_2}}

\end{ELimportant}

\end{ELjoin}

\begin{ELunit}

\begin{itemize}
\tightlist
\item
  When a steady current crosses the boundary between two different lossy dielectrics
\end{itemize}

\ELdisp{J_{1n}\ELbr =J_{2n}\to\sigma_1E_{1n}\ELbr =\sigma_2E_{2n}}
\ELdisp{D_{1n}-D_{2n}\ELbr =\rho_s\to\epsilon_1E_{1n}-\epsilon_2E_{2n}\ELbr =\rho_s}
\ELdisp{\rho_s\ELbr =\left(\epsilon_1\frac{\sigma_2}{\sigma_1}-\epsilon_2\right)E_{2n}\ELbr =\left(\epsilon_1-\epsilon_2\frac{\sigma_1}{\sigma_2}\right)E_{1n}}

\begin{itemize}
\tightlist
\item
  \ELim{E_{1n}=\uvec{n2}\cdot\vect{E}_1} and \ELim{E_{2n}=\uvec{n2}\cdot\vect{E}_2}
\item
  If medium 2 is a much better conductor than medium 1
\end{itemize}

\ELdisp{\rho_s\ELbr =\epsilon_1E_{1n}\ELbr =D_{1n}}

\end{ELunit}

\begin{ELjoin}

\begin{itemize}
\tightlist
\item
  which is the boundary condition of a conductor.
\end{itemize}

\ELnextnum{4-1}

\begin{ELexample}{}

A potential difference \ELim{\mathcal{V}} is applied to a parallel-plate capacitor of area \ELim{S}. The space between the conducting plates is filled with two lossy dielectric materials of thicknesses \ELim{d_1} and \ELim{d_2}, permittivities \ELim{\epsilon_1} and \ELim{\epsilon_2}, and conductivities \ELim{\sigma_1} and \ELim{\sigma_2}. Find

\begin{itemize}
\item
  \begin{enumerate}
  \def\labelenumi{(\alph{enumi})}
  \tightlist
  \item
    the electric field intensity in the two media
  \end{enumerate}
\item
  \begin{enumerate}
  \def\labelenumi{(\alph{enumi})}
  \setcounter{enumi}{1}
  \tightlist
  \item
    the current density in the two media
  \end{enumerate}
\item
  \begin{enumerate}
  \def\labelenumi{(\alph{enumi})}
  \setcounter{enumi}{2}
  \tightlist
  \item
    the surface charge densities on the conducting plates and at the interface of the two dielectrics
  \end{enumerate}
\end{itemize}

\ELfigure{m04-ex1}{}

\begin{ELsolution}

\begin{itemize}
\item
  \begin{enumerate}
  \def\labelenumi{(\alph{enumi})}
  \tightlist
  \item
  \end{enumerate}
\end{itemize}

\ELdisp{J_1\ELbr =J_2\quad\ELbr \ELbr \Longrightarrow\quad\ELbr \begin{aligned}\mathcal{V}&=E_1d_1+E_2d_2\\\sigma_1E_1&=\sigma_2E_2\end{aligned}}
\ELdisp{E_1\ELbr =\frac{\sigma_2\mathcal{V}}{\sigma_2d_1+\sigma_1d_2}\qquad\ELbr  E_2\ELbr =\frac{\sigma_1\mathcal{V}}{\sigma_2d_1+\sigma_1d_2}}

\begin{itemize}
\item
  \begin{enumerate}
  \def\labelenumi{(\alph{enumi})}
  \setcounter{enumi}{1}
  \tightlist
  \item
  \end{enumerate}
\end{itemize}

\ELdisp{J_1\ELbr =J_2\ELbr =\sigma_1E_1\ELbr =\sigma_2E_2\ELbr =\frac{\sigma_1\sigma_2\mathcal{V}}{\sigma_2d_1+\sigma_1d_2}}

\begin{itemize}
\item
  \begin{enumerate}
  \def\labelenumi{(\alph{enumi})}
  \setcounter{enumi}{2}
  \tightlist
  \item
  \end{enumerate}
\end{itemize}

\ELdisp{\uvec{n}\cdot\vect{E}_1\ELbr =\frac{\rho_{s1}}{\epsilon_1}\quad\ELbr \ELbr \Longrightarrow\quad\ELbr \rho_{s1}\ELbr =\epsilon_1E_1\ELbr =\frac{\epsilon_1\sigma_2\mathcal{V}}{\sigma_2d_1+\sigma_1d_2}}
\ELdisp{\uvec{n}\cdot\vect{E}_2\ELbr =\frac{\rho_{s2}}{\epsilon_2}\quad\ELbr \ELbr \Longrightarrow\quad\ELbr \rho_{s2}\ELbr =-\epsilon_2E_2\ELbr =-\frac{\epsilon_2\sigma_1\mathcal{V}}{\sigma_2d_1+\sigma_1d_2}}
\ELdisp{\rho_{si}\ELbr =\left(\epsilon_1-\epsilon_2\frac{\sigma_1}{\sigma_2}\right)E_{1n}\quad\ELbr \ELbr \Longrightarrow\quad\ELbr \begin{aligned}\rho_{si}&=\left(\epsilon_2\frac{\sigma_1}{\sigma_2}-\epsilon_1\right)\frac{\sigma_2\mathcal{V}}{\sigma_2d_1+\sigma_1d_2}\\&=\frac{(\epsilon_2\sigma_1-\epsilon_1\sigma_2)\mathcal{V}}{\sigma_2d_1+\sigma_1d_2}\qquad(\mathrm{C/m^2})\end{aligned}}

\end{ELsolution}

\end{ELexample}

\end{ELjoin}

\ELhold

\ELsection{Resistance Calculations}{Resistance Calculations}

\begin{ELunit}

\ELfigure{m04-map-resistance}{}

\begin{itemize}
\item
  How is the \textbf{leakage resistance} of capacitive structures computed?
\item
  How is the resistance of various structures computed?
\item
  The capacitance between two conductors separated by a dielectric medium is
\end{itemize}

\ELdisp{C\ELbr =\frac QV\ELbr =\frac{\oint_S\vect{D}\cdot d\vect{s}}{-\int_L\vect{E}\cdot\dif\vect{\ell}}\ELbr =\frac{\oint_S\epsilon\vect{E}\cdot d\vect{s}}{-\int_L\vect{E}\cdot\dif\vect{\ell}}}

\begin{itemize}
\tightlist
\item
  The surface integral is taken over a surface enclosing the positive conductor, and the line integral from the negative conductor to the positive conductor.
\item
  If the dielectric medium is lossy, a current flows from the positive to the negative conductor. The resistance between the conductors (the leakage resistance) is
\end{itemize}

\ELdisp{R\ELbr =\frac VI\ELbr =\frac{-\int_L\vect{E}\cdot\dif\vect{\ell}}{\oint_S\vect{J}\cdot d\vect{s}}\ELbr =\frac{-\int_L\vect{E}\cdot\dif\vect{\ell}}{\oint_S\sigma\vect{E}\cdot d\vect{s}}}

\end{ELunit}

\begin{ELjoin}

\begin{itemize}
\tightlist
\item
  If the permittivity and the conductivity have the same space dependence (\ELim{\epsilon(x,y,z)=k\sigma(x,y,z)}), or if the dielectric medium is homogeneous (\ELim{\epsilon} and \ELim{\sigma} independent of the space coordinates)
\end{itemize}

\begin{ELimportant}{}

\ELdisp{RC\ELbr =\frac CG\ELbr =\frac\epsilon\sigma}

\end{ELimportant}

\end{ELjoin}

\ELnextnum{4-2}

\begin{ELexample}{}

Find the leakage resistance per unit length

\begin{itemize}
\item
  \begin{enumerate}
  \def\labelenumi{(\alph{enumi})}
  \tightlist
  \item
    between the inner and outer conductors of a coaxial cable with inner-conductor radius \ELim{a}, outer-conductor radius \ELim{b} and a dielectric of conductivity \ELim{\sigma}
  \end{enumerate}
\item
  \begin{enumerate}
  \def\labelenumi{(\alph{enumi})}
  \setcounter{enumi}{1}
  \tightlist
  \item
    of a two-wire transmission line consisting of two wires of radius \ELim{a} separated by \ELim{D}, in a medium of conductivity \ELim{\sigma}
  \end{enumerate}
\end{itemize}

\begin{ELsolution}

\begin{itemize}
\item
  \begin{enumerate}
  \def\labelenumi{(\alph{enumi})}
  \tightlist
  \item
    The capacitance per unit length of a coaxial cable
  \end{enumerate}
\end{itemize}

\ELdisp{C_1\ELbr =\frac{2\pi\epsilon}{\ln(b/a)}\qquad\ELbr (\mathrm{F/m})}

\begin{itemize}
\tightlist
\item
  The leakage resistance per unit length
\end{itemize}

\ELdisp{R_1\ELbr =\frac\epsilon\sigma\left(\frac{1}{C_1}\right)\ELbr =\frac{1}{2\pi\sigma}\ln\left(\frac ba\right)\qquad\ELbr (\Omega\cdot\mathrm{m})}

\begin{itemize}
\item
  \begin{enumerate}
  \def\labelenumi{(\alph{enumi})}
  \setcounter{enumi}{1}
  \tightlist
  \item
    The capacitance per unit length of a two-wire transmission line
  \end{enumerate}
\end{itemize}

\ELdisp{C'_1\ELbr =\frac{\pi\epsilon}{\cosh^{-1}\left(\dfrac{D}{2a}\right)}\qquad\ELbr (\mathrm{F/m})}

\begin{itemize}
\tightlist
\item
  The leakage resistance per unit length
\end{itemize}

\ELdisp{\begin{aligned}R'_1&=\frac\epsilon\sigma\left(\frac{1}{C'_1}\right)=\frac{1}{\pi\sigma}\cosh^{-1}\left(\frac{D}{2a}\right)\\&=\frac{1}{\pi\sigma}\ln\left[\frac{D}{2a}+\sqrt{\left(\frac{D}{2a}\right)^2-1}\right]\qquad(\Omega\cdot\mathrm{m})\end{aligned}}

\end{ELsolution}

\end{ELexample}

\begin{ELunit}

\begin{itemize}
\tightlist
\item
  In certain situations electrostatic and steady-current problems are not exactly analogous.
\item
  \textbf{First method} for computing the resistance of a piece of conducting material between two surfaces or terminals

  \begin{itemize}
  \tightlist
  \item
    Choosing an appropriate coordinate system
  \item
    Applying a potential difference \ELim{V_0} between the two terminals
  \item
    Finding the electric field intensity (if the material is homogeneous, by solving Laplace's equation for the electric potential and then computing the electric field intensity)
  \item
    Finding the total current
    \ELdisp{I\ELbr =\int_S\vect{J}\cdot d\vect{s}\ELbr =\int_S\sigma\vect{E}\cdot d\vect{s}}
  \item
    \ELim{S}: the cross-section through which the current flows
  \item
    The resistance is \ELim{V_0/I}
  \end{itemize}
\item
  \textbf{Second method} for computing the resistance of a piece of conducting material between two surfaces or terminals (when \ELim{J} can easily be determined from \ELim{I})

  \begin{itemize}
  \tightlist
  \item
    Choosing an appropriate coordinate system
  \item
    Passing a current \ELim{I} between the two terminals
  \item
    Finding \ELim{J} from \ELim{I}
  \item
    Computing the electric field intensity \ELim{\vect{E}=\vect{J}/\sigma} and the potential difference \ELim{V_0}
    \ELdisp{V_0\ELbr =-\int\vect{E}\cdot\dif\vect{\ell}}

    \begin{itemize}
    \tightlist
    \item
      The integration is carried out from the terminal at the lower potential to the terminal at the higher potential.
    \end{itemize}
  \item
    The resistance is \ELim{V_0/I}
  \end{itemize}
\end{itemize}

\end{ELunit}

\ELnextnum{4-3}

\begin{ELexample}{}

Consider a conducting material of uniform thickness \ELim{h} and conductivity \ELim{\sigma} in the shape shown in the figure below. Find the resistance between its two end faces.

\ELfigure{m04-ex3}{}

\begin{ELsolution}

\begin{itemize}
\tightlist
\item
  Choosing cylindrical coordinates
\end{itemize}

\ELdisp{V\ELbr =0\quad\ELbr \text{at}\quad\ELbr \phi\ELbr =0,\qquad\ELbr  V\ELbr =V_0\quad\ELbr \text{at}\quad\ELbr \phi\ELbr =\pi/2}
\ELdisp{\frac{d^2V}{d\phi^2}\ELbr =0\quad\ELbr \ELbr \Longrightarrow\quad\ELbr  V\ELbr =c_1\phi+c_2\quad\ELbr \ELbr \Longrightarrow\quad\ELbr  V\ELbr =\frac{2V_0}{\pi}\phi}
\ELdisp{\begin{aligned}\vect{J}=\sigma\vect{E}&=-\sigma\nabla V\\&=-\uvec{\phi}\sigma\frac{\partial V}{r\,\partial\phi}=-\uvec{\phi}\frac{2\sigma V_0}{\pi r}\end{aligned}}
\ELdisp{\begin{aligned}I=\int_S\vect{J}\cdot d\vect{s}&=\int_0^h\!\!\int_a^b\left(-\uvec{\phi}\frac{2\sigma V_0}{\pi r}\right)\cdot\left(-\uvec{\phi}\,dr\,dz\right)\\&=\frac{2\sigma hV_0}{\pi}\ln\frac ba\end{aligned}}
\ELdisp{R\ELbr =\frac{V_0}{I}\ELbr =\frac{\pi}{2\sigma h\ln(b/a)}}

\end{ELsolution}

\end{ELexample}
